\section*{\texorpdfstring{\S\CurrentExerciseGroup}{第{\CurrentExerciseGroup}节}}
\phantomsection
\addcontentsline{toc}{section}{第{\CurrentExerciseGroup}节习题}

\begin{enumerate}
\def\labelenumi{\arabic{enumi})}
\item
  考察与区间 {[}0, 1{]} 相同的局部紧空间 T、T'，以及 T 与 T' 上分别与 Lebesgue 测度相同的测度 \(\mu, \mu'\)。设 F 是一个 Hilbert 空间，具有可数标准正交基，并将其排成双重序列 \((\mathbf{e}_{mn})\)；设 F' = F 为其对偶。令 \(\mathbf{u}_m(t) = \mathbf{e}_{mn}\)，适用于 \((n-1)2^{-m} \leqslant t < n2^{-m}\) 与 \(1 \leqslant n \leqslant 2^m\)；令 \(\mathbf{f}(t, t') = 2^m \mathbf{u}_m(t)\)，适用于 \(2^{-m} < t' \leqslant 2^{-m+1}\)，最后令 \(\mathbf{f}(t, 0) = 0\) 与 \(\mathbf{f}(1, t') = 0\)。试证 \(\mathbf{f}\) 关于乘积测度 \(\mu \otimes \mu'\) 是标量可积的，但函数 \(t' \mapsto \mathbf{f}(t, t')\) 不是标量 \(\mu'\)-可积的——对任何 \(t \in T\) 都是如此。
\item
  设 \(\mathrm{T}, \mathrm{X}, \mathrm{Y}\) 是三个局部紧空间，其中 \(\mathrm{Y}\) 具有可数基。设 \(\mu\) 是一个测度 \(\geqslant 0\)，定义在 \(\mathrm{T}\) 上；设 \(t \mapsto \lambda_t (t \in \mathrm{T})\) 是一个 \(\mu\)-适宜\(^1\) 的测度族，\(\geqslant 0\)，定义在 \(\mathrm{X}\) 上，并令 \(\nu = \int \lambda_t d\mu(t)\)。设 \(x \mapsto \rho_x (x \in \mathrm{X})\) 是一个 \(\nu\)-适宜的测度族，\(\geqslant 0\)，定义在 \(\mathrm{Y}\) 上。此外，假设还满足下列两个条件之一：a) \(\mathrm{X}\) 是无穷远处可数的；
\end{enumerate}

\begin{enumerate}
\def\labelenumi{\alph{enumi})}
\setcounter{enumi}{1}
\tightlist
\item
  \(\nu\) 是有界的。
\end{enumerate}

在这些条件下，试证存在一个局部 \(\mu\)-可忽略集 \(N'\)，使得对每个 \(t \notin N'\)，族 \(x \mapsto \rho_x\) 都是 \(\lambda_t\)-适宜的；此外，族 \(t \mapsto \int \rho_x d\lambda_t(x)\)（对 \(t \notin N'\) 有定义）是 \(\mu\)-适宜的，且

\[
\int d \mu (t) \int \rho_ {x} d \lambda_ {t} (x) = \int \rho_ {x} d \nu (x).
\]

（为证 \(x \mapsto \rho_x\) 局部几乎处处为标量 \(\lambda_t\)-可积，可利用 §3 第 1 号的引理 1，在情形 \(b\) 中要利用 \(\lambda_t\) 局部几乎处处有界这一事实；用同样的方法可证 \(t \mapsto \int \rho_x d\lambda_t(x)\) 是标量 \(\mu\)-可积的；再利用第 5 号的命题 13。）

\begin{enumerate}
\def\labelenumi{\arabic{enumi})}
\setcounter{enumi}{2}
\tightlist
\item
  设 S、T、Y 是三个无穷远处可数的局部紧空间，且 Y 具有可数基。设 \(\rho\)（相应地 \(\sigma\)）是 S（相应地 T）上的正测度，\(\nu = \rho \otimes \sigma\)，并设 \((\lambda_{s,t})_{(s,t) \in \mathbb{S} \times \mathbb{T}}\) 是 Y 上的一个 \(\nu\)-适宜\(^{1}\)正测度族。那么，存在一个 \(\rho\)-可忽略集 N，使得对每个 \(s \notin N\)，族 \((\lambda_{s,t})_{t \in \mathbb{T}}\) 都是 \(\sigma\)-适宜的；族 \(s \mapsto \int \lambda_{s,t} d\sigma(t)\)（几乎处处有定义）是 \(\rho\)-适宜的，且
\end{enumerate}

\[
\iint \lambda_ {s, t} d \rho (s) d \sigma (t) = \int d \rho (s) \int \lambda_ {s, t} d \sigma (t).
\]

（对空间 \(\mathrm{X} = \mathrm{S} \times \mathrm{T}\) 应用习题 2，并注意 \(\nu = \int (\varepsilon_s \otimes \sigma) d\rho(s)\)。）

\begin{enumerate}
\def\labelenumi{\arabic{enumi})}
\setcounter{enumi}{3}
\tightlist
\item
  设 \(\mathrm{T}, \mathrm{X}, \mathrm{Y}\) 是三个局部紧空间，其中 \(\mathrm{X}\) 是无穷远处可数的。设 \(\mu\) 是 \(\mathrm{T}\) 上的正测度，\((\lambda_t)_{t \in \mathrm{T}}\) 是一个 \(\mu\)-适宜正测度族（在 \(\mathrm{X}\) 上），且 \(\nu = \int \lambda_t d\mu(t)\)。假设下列两个条件之一成立：
\end{enumerate}

\begin{enumerate}
\def\labelenumi{\alph{enumi})}
\item
  Y 具有可数基；
\item
  测度 \(\nu\) 是有界的。
\end{enumerate}

在这些条件下，若 \(\pi\) 是 X 到 Y 中的 \(\nu\)-逆紧映射，则诸 \(t \in \mathbf{T}\) 中使 \(\pi\) 不是 \(\lambda_t\)-逆紧者所成之集是局部 \(\mu\)-可忽略的；此外，Y 上的正测度族 \((\pi(\lambda_t))_{t \in \mathbf{T}}\)（在 T 中局部几乎处处有定义）是 \(\mu\)-适宜的，且

\[
\pi \left(\int \lambda_ {t} d \mu (t)\right) = \int \pi (\lambda_ {t}) d \mu (t).
\]

（在情形 \(a\)，对 \(\rho_x = \varepsilon_{\pi(x)}\) 应用习题 2。在情形 \(b\)，注意 \(\lambda_t\) 局部几乎处处有界，于是可归结为证明映射 \(t \mapsto \pi(\lambda_t)\) 是淡 \(\mu\)-可测的（参见第 V 章 §3 第 2 号命题 4）。给定一个数 \(\varepsilon > 0\)，并设 \(K\) 是 \(T\) 的紧子集，首先注意 \(X\) 中存在递增的紧集序列 \((F_m)\)，使 \(\pi\) 在每个 \(F_m\) 上的限制连续，且令 \(N_m = X - F_m\) 时有 \(\nu(N_m) \leqslant \varepsilon/4^m\)。设 \(A_m\) 是 \(t \in K\) 中使 \(N_m\) 不是 \(\lambda_t\)-可积或使 \(\lambda_t(N_m) > 1/2^m\) 的那些 t 所成的集；若 \(A\) 是诸 \(A_m\) 之并，则 \(\mu(A) \leqslant \varepsilon/2\)。设 \(K_1\) 是 \(K - A\) 的紧子集，使 \(\mu(K - K_1) \leqslant \varepsilon\)，且对每个函数 \(g \in \mathcal{H}(X)\)，\(t \mapsto \int g(x) d\lambda_t(x)\) 在 \(K_1\) 上的限制连续。利用 Urysohn 定理，试证对每个函数 \(f \in \mathcal{H}(Y)\)，在 \(K_1\) 上 \(t \mapsto \int f(\pi(x)) d\lambda_t(x)\) 的限制是连续函数的一致极限。）

*5) 对每个 \(t \in \mathbf{R}\)，令 \(\mathbf{f}(t)\) 为连续函数 \(x \mapsto \exp(-2\pi itx)\)，它是空间 \(\mathcal{C} = \mathcal{C}_{\mathbf{C}}(\mathbf{R})\)（\(\mathbf{R}\) 上的连续复函数空间）中的一个元素。空间 \(\mathcal{C}\) 与空间 \(\mathcal{C}' = \mathcal{C}'_{\mathbf{C}}(\mathbf{R})\)（\(\mathbf{R}\) 上具有紧支集的复测度组成的空间）成对偶；具有紧支集的连续复函数空间 \(\mathcal{K} = \mathcal{K}_{\mathbf{C}}(\mathbf{R})\) 可按典范方式与 \(\mathcal{C}'\) 的一个子空间等同，办法是把函数 \(\varphi \in \mathcal{K}\) 等同于关于 Lebesgue 测度以 \(\varphi\) 为密度的测度。用 \(\mathcal{D}\) 表示由 \(\mathcal{K}\) 中具有紧支集的无穷次可微复函数所构成的子空间。试证：当 \(\mathcal{C}\) 赋予拓扑 \(\sigma(\mathcal{C},\mathcal{C}')\) 或拓扑 \(\sigma(\mathcal{C},\mathcal{K})\) 时，\(\mathbf{f}\) 关于 Lebesgue 测度 \(\mu\) 不是标量可积的；然而对拓扑 \(\sigma(\mathcal{C},\mathcal{D})\)，\(\mathbf{f}\) 是标量 \(\mu\)-可积的，且 \(\int f d\mu\) 是测度 \(\varepsilon_0\)；试证在此情形命题 7 的条件得到验证。*

\begin{enumerate}
\def\labelenumi{\arabic{enumi})}
\setcounter{enumi}{5}
\item
  设 F 是一个可分辨的 Fréchet 空间（TVS，IV，§2，习题 4 b）），F' 是其对偶。试证：若 f 是 T 到 F 中的标量本质可积映射，则 \(\int f d\mu\) 属于 F 的双对偶 F'\,'。（把 F 嵌入 F'\,'；应用定理 1 以及命题 2 与附录的引理。）
\item
  \begin{enumerate}
  \def\labelenumii{\alph{enumii})}
  \tightlist
  \item
    设 \(\mathbf{F} = \overline{\mathcal{K}(\mathbf{N})}\) 是趋于 0 的实数序列组成的 Banach 空间，\(\mathbf{F}' = \mathbf{L}^1 (\mathbf{N})\) 是其对偶（TVS，IV，§1，习题 1）。设 \(\mathbf{f} = (f_n)\) 是 \(\mathrm{I} = [0,1]\) 到 \(\mathbf{F}\) 中的映射，满足 \(f_{n}(t) = n\varphi_{\mathrm{I}_{n}}(t)\)，其中 \(\mathrm{I}_n = [0,1 / n]\)。试证 \(\mathbf{f}\) 关于 Lebesgue 测度 \(\mu\)（在 \(\mathrm{I}\) 上）是标量可积的，但 \(\int \mathbf{f}d\mu\) 是 \(\mathbf{F}''\) 中不属于 \(\mathbf{F}\) 的元素。
  \end{enumerate}
\end{enumerate}

\begin{enumerate}
\def\labelenumi{\alph{enumi})}
\setcounter{enumi}{1}
\item
  设 \(\mathbf{g}\) 是 \(J = [-1, +1]\) 到 \(F\) 中的映射，由条件 \(\mathbf{g}(t) = \mathbf{f}(t)\)（对 \(t \geqslant 0\)）与 \(\mathbf{g}(t) = -\mathbf{f}(-t)\)（对 \(t \leqslant 0\)）定义。则 \(\mathbf{g}\) 关于 \(J\) 上的 Lebesgue 测度是标量可积的，且 \(\int \mathbf{g} d\mu \in F\)，但存在函数 \(h \in \mathcal{L}^{\infty}\) 使 \(\int h\mathbf{g} d\mu \notin F\)。
\item
  由 \(a\) 推出拓扑向量空间 \(F\) 不与任何赋范空间的对偶同构这一事实的一个新证明（TVS，IV，§2，习题 22 c））。
\end{enumerate}

\begin{enumerate}
\def\labelenumi{\arabic{enumi})}
\setcounter{enumi}{7}
\tightlist
\item
  设 \(\mathbf{F}\) 是一个 Hausdorff 局部凸空间，\(\mathbf{f}\) 是 \(\mathbf{T}\) 到 \(\mathbf{F}\) 中的标量局部可积映射，使得对每个紧子集 \(\mathbf{K}\)（\(\mathbf{T}\) 中的）都有 \(\int_{\mathbf{K}} \mathbf{f} d\mu \in \mathbf{F}\)。
\end{enumerate}

\begin{enumerate}
\def\labelenumi{\alph{enumi})}
\item
  试证：若 \(\mathbf{f}\) 是标量本质可积的，且 \(\mathrm{F}\) 是半自反的（或者只需每个对 \(\sigma (\mathrm{F},\mathrm{F}^{\prime})\) 的 Cauchy 序列对该拓扑都在 \(\mathrm{F}\) 中收敛，当 \(\mathrm{T}\) 是无穷远处可数时），则 \(\int g\mathbf{f}d\mu \in \mathrm{F}\) 对每个函数 \(g\in \mathcal{L}^{\infty}\) 成立。
\item
  试证：若 F 是拟完备的，且对 F 上每个连续半范数 q 都有 \(\int^{\bullet} q(\mathbf{f}(t)) d\mu(t) < +\infty\)，则 f 是标量本质可积的，且 \(\int g\mathbf{f} d\mu \in F\) 对每个函数 \(g \in L^{\infty}\) 成立。
\end{enumerate}

（在两种情形中，都考察 T 的紧子集组成的递增有向集。）

\begin{enumerate}
\def\labelenumi{\arabic{enumi})}
\setcounter{enumi}{8}
\item
  设 G、H 是两个 Hausdorff 局部凸空间，U 是 T 到 \(F = \mathcal{L}_{s}(G;H)\) 中的映射。假设 G 是桶式的，H 是拟完备的，U 是 \(\mu\)-可测的，且对 T 的每个紧子集 K，\(U(K)\) 都有界；此外还假设，对 H 上每个连续半范数 q 以及每个 \(x \in G\)，\(\int^{\bullet} q(U(t) \cdot x) d\mu(t) < +\infty\)。在这些条件下，试证 U 是标量本质 \(\mu\)-可积的，且 \(\int U(t) d\mu(t) \in F\)（利用习题 8 b））。
\item
  设 \(G, H\) 是两个 Fréchet 空间，\(G_0\)（相应地 \(H_0\)）是 \(G\)（相应地 \(H\)）的稠密子集，\(t \mapsto \Phi_t\) 是 \(T\) 到空间 \(F = \mathcal{B}(G, H)\)（即 \(G \times H\) 上连续双线性型组成的空间，赋以逐点收敛拓扑）中的映射。假设：\(1^\circ\) 对每个对 \((a, b) \in G_0 \times H_0\)，\(t \mapsto \Phi_t(a, b)\) 都是本质 \(\mu\)-可积的；\(2^\circ\) 若对每对有界子集 \(A \subset G\)、\(B \subset H\)，令 \(q_{A,B}(\Phi) = \sup_{(x, y) \in A \times B} |\Phi(x, y)|\)，其中
\end{enumerate}

对每个 \(\Phi \in \mathcal{B}(\mathrm{G},\mathrm{H})\) 都有 \(\int^{*}q_{\mathrm{A},\mathrm{B}}(\Phi_t)d\mu (t) < + \infty\)。在这些条件下，试证 \(t\mapsto \Phi_t\) 是标量本质 \(\mu\)-可积的，且 \(\int \Phi_t d\mu (t)\in \mathrm{F}\)（利用 Lebesgue 定理，归结为应用定理 1 的推论 1）。\(\mathrm{H} = \mathbf{R}\) 的特殊情形。

\begin{enumerate}
\def\labelenumi{\arabic{enumi})}
\setcounter{enumi}{10}
\item
  设 \(\mu\) 是 \(T = [0,1]\) 上的 Lebesgue 测度，\(F\) 是具有可数标准正交基 \((\mathbf{e}_n)\) 的 Hilbert 空间。设 \(\mathbf{f}\) 是 \(T\) 到 \(F\) 中的映射，满足 \(\mathbf{f}(0) = 0\)，且 \(\mathbf{f}(t) = 2^n \mathbf{e}_n / n\) 在区间 \([2^{-n-1}, 2^{-n}]\) 上成立（\(n \geqslant 0\)）。试证 \(\mathbf{f}\) 是 \(\mu\)-可测的、标量 \(\mu\)-可积的，且 \(\int \mathbf{f} d\mu \in F\)，但 \(\int^* |\mathbf{f}| d\mu = +\infty\)。
\item
  设 \(\mu\) 是 \(T = [0,1]\) 上的 Lebesgue 测度，并设 \(F\) 是一个 Hilbert 空间，具有标准正交基 \((\mathbf{e}_t)_{t \in T}\)，与 \(T\) 等势。
\end{enumerate}

\begin{enumerate}
\def\labelenumi{\alph{enumi})}
\item
  设 \(\mathbf{f}\) 是映射 \(t \mapsto \mathbf{e}_t\)，从 \(\mathrm{T}\) 到 \(\mathrm{F}\) 中。试证 \(\mathbf{f}\) 是标量 \(\mu\)-可忽略的，但不是 \(\mu\)-可测的（当 \(\mathrm{F}\) 赋以弱拓扑 \(\sigma(\mathrm{F}, \mathrm{F}')\) 时；更不用说当 \(\mathrm{F}\) 赋以其原拓扑时了）。
\item
  设 A 是 T 中的一个不可测集（第 IV 章 §4，习题 8）；试证函数 \(\mathbf{g} = \mathbf{f}\varphi_{\mathrm{A}}\) 是标量 \(\mu\)-可忽略的，但数值函数 \(|\mathbf{g}|\) 不是 \(\mu\)-可测的。
\end{enumerate}

\begin{enumerate}
\def\labelenumi{\arabic{enumi})}
\setcounter{enumi}{12}
\item
  设 F 是 Hausdorff 局部凸空间，\(F'\) 是其对偶，q 是 F 上的一个下半连续半范数。试证：若 f 是 T 到 F 中的映射，是 \(\mu\)-可测的（对拓扑 \(\sigma(\mathrm{F},\mathrm{F}')\)），则数值函数 \(q \circ f\) 是 \(\mu\)-可测的。
\item
  设 \(\mu\) 是 T = {[}0,1{]} 上的 Lebesgue 测度；用 F 表示由 T 上有限的 \(\mu\)-可测数值函数组成的 R 上的向量空间，赋以逐点收敛拓扑，这使它成为一个 Hausdorff 局部凸空间。
\end{enumerate}

\begin{enumerate}
\def\labelenumi{\alph{enumi})}
\item
  试证 F 中存在可数稠密子集（考察 T 上连续数值函数组成的 Banach 空间 \(\mathcal{C}(T)\) 中的一个稠密序列）。
\item
  对每个 \(t \in T\)，令 \(\mathbf{f}(t)\) 是 F 的对偶 \(F'\) 中由 \(\langle\mathbf{f}(t), z\rangle = z(t)\)（对一切 \(t \in T\)）定义的元素。试证：当 \(F'\) 赋以拓扑 \(\sigma(\mathrm{F}', \mathrm{F})\) 时，f 是标量 \(\mu\)-可测的，但不是 \(\mu\)-可测的（参见命题 13）。
\end{enumerate}

\begin{enumerate}
\def\labelenumi{\arabic{enumi})}
\setcounter{enumi}{14}
\item
  设 F 是可度量化的局部凸空间，f 是 T 到 F 中的标量 \(\mu\)-可测映射，使得对 T 的每个紧子集 K，存在 F 的可数子集 H，使 \(\mathbf{f}(t) \in \overline{\mathbf{H}}\) 对几乎所有的 \(t \in K\) 成立；试证在这些条件下，f 对 F 的原拓扑是 \(\mu\)-可测的。（在化归到 F 为可分的情形之后，把 F 嵌入可数个赋范空间的乘积中。）
\item
  设 \(\mathbf{F}\) 是 Banach 空间，\(\mathbf{F}'\) 是其对偶。对每个 \(p\)，当 \(1 \leqslant p \leqslant +\infty\) 时，用 \(\Lambda_{\mathbf{F}'}^p(\mathrm{T},\mu)\)（或简记为 \(\Lambda_{\mathbf{F}'}^p\)）表示由这样的映射 \(\mathbf{f}\)——即 \(\mathrm{T}\) 到 \(\mathbf{F}'\) 中的、使得对每个 \(\mathbf{z} \in \mathbf{F}\) 数值函数 \(\langle \mathbf{z},\mathbf{f}\rangle\) 属于 \(\mathcal{L}^p(\mathrm{T},\mu)\) 的映射——组成的集。有 \(\Lambda_{\mathbf{F}'}^p \supset \mathcal{L}_{\mathbf{F}'}^p\)，且 \(\Lambda_{\mathbf{F}'}^1\) 是标量本质 \(\mu\)-可积的、取值于 \(\mathbf{F}'\) 的函数组成的空间（\(\mathbf{F}'\) 赋以 \(\sigma(\mathbf{F}',\mathbf{F})\)）。已知，若用 \(\theta(\mathbf{z})\) 表示 \(\langle \mathbf{z},\mathbf{f}\rangle\) 在 \(\mathrm{L}^p(\mu)\) 中的类，则映射 \(\mathbf{z} \mapsto \theta(\mathbf{z})\) 是连续的（第 4 号，引理 2）；设 \(M_p(\mathbf{f})\) 为其范数；它是空间 \(\Lambda_{\mathbf{F}'}^p\) 上的一个半范数。为使 \(M_p(\mathbf{f}) = 0\)，必要且充分的是 \(\mathbf{f}\) 为标量局部可忽略的。
\end{enumerate}

\begin{enumerate}
\def\labelenumi{\alph{enumi})}
\item
  为使 \(\mathbf{f} \in \Lambda_{\mathbf{F}'}^p\)，必要且充分的是：对每个数值函数 \(g \in \mathcal{L}^q\)，函数 \(g\mathbf{f}\) 都是标量本质 \(\mu\)-可积的；此时有 \(\mathrm{M}_1(g\mathbf{f}) \leqslant \mathrm{M}_p(\mathbf{f})\mathrm{N}_q(g)\)。
\item
  在空间 \(\Lambda_{\mathbf{F}'}^1\) 中，试证半范数 \(\mathbf{M}_1\) 等价于半范数 \(\mathrm{M}_1'(\mathbf{f}) = \sup |\int_{\mathbf{A}} \mathbf{f} d\mu|\)，其中 A 取遍 T 的可测子集所成的集；更精确地，\(\mathrm{M}_1' \leqslant \mathrm{M}_1 \leqslant 2\mathrm{M}_1'\)。由此推出 \(\mathrm{M}_p\) 是等价于 \(\mathrm{M}_p'(\mathbf{f}) = \sup_{\mathrm{N}_q(g) \leqslant 1} \mathrm{M}_1'(g\mathbf{f})\) 的半范数。
\item
  取 F 为一个具有可数标准正交基的 Hilbert 空间，取 \(\mu\) 为 \(T = [0,1]\) 上的 Lebesgue 测度。试证存在 T 到 F 中的一个映射，它不可积但标量可积，并且是可测阶梯函数序列在半范数 \(M_1\) 下的极限（参见习题 11）；由此推出：在空间 \(\mathcal{L}_{\mathrm{F}}^1\)（即 \(\mu\)-可积函数组成的空间）上，由半范数 \(M_1\) 定义的拓扑严格粗于由 \(N_1\) 定义的拓扑。对半范数 \(M_p\) 与 \(N sub p\)（\(1 \leqslant p < +\infty\)）有类似的结果。
\end{enumerate}

\(d\)）试证：在空间 \(\mathcal{L}_{\mathrm{F}^{\prime}}^{\infty}\) 上，\(\mathrm{M}_{\infty} = \mathrm{N}_{\infty}\)

\begin{enumerate}
\def\labelenumi{\alph{enumi})}
\setcounter{enumi}{4}
\tightlist
\item
  取 F 为一个具有可数标准正交基的 Hilbert 空间，该基排成双重序列 \((\mathbf{e}_{mn})\)。设 \(\mu\) 是 \(T = [0,1]\) 上的 Lebesgue 测度。设 \(\mathbf{u}_m\) 是 T 到 F 中的映射，满足 \(\mathbf{u}_m(t) = \mathbf{e}_{mn}\)（对 \((n - 1)2^{-m} \leqslant t < n2^{-m}\) 及 \(1 \leqslant n \leqslant 2^m\)），且 \(\mathbf{u}_m(1) = 0\)；令 \(\mathbf{f}_n = \sum_{i=0}^{n} \mathbf{u}_i\)。试证：对 \(1 \leqslant p < +\infty\)，\((\mathbf{f}_n)\) 是 \(\Lambda_{\mathrm{F}'}^p\) 中的 Cauchy 序列，但不收敛于 \(\Lambda_{\mathrm{F}'}^p\) 中的任何函数。
\end{enumerate}

\begin{enumerate}
\def\labelenumi{\alph{enumi})}
\setcounter{enumi}{5}
\item
  设 \((\mathbf{f}_{n})\) 是空间 \(\Lambda_{F'}^{p}\)（\(1 \leqslant p \leqslant +\infty\)）中的 Cauchy 序列；假设存在 T 到 \(F'\) 中的映射 f，使得对每个 \(z \in F\)，函数序列 \(\langle z, f_{n} \rangle\) 都依测度收敛于 \(\langle z, f \rangle\)（第 IV 章 §5 第 11 号）。试证 \(f \in \Lambda_{F'}^{p}\)，且序列 \((\mathbf{f}_{n})\) 在 \(\Lambda_{F'}^{p}\) 中以 f 为极限。
\item
  取 F 为绝对收敛级数组成的 Banach 空间 \(\mathrm{L}^{1}(\mathbf{N})\)；其对偶为 \(\mathbf{F}^{\prime} = \mathrm{L}^{\infty}(\mathbf{N})\)（TVS，IV，§1，习题 1）。设 \(\mu\) 是 T = {[}0,1{]} 上的 Lebesgue 测度。设 f 是 T 到 \(F^{\prime}\) 中的映射，满足 \(\mathbf{f}(0) = 0\)，且对 \(2^{-n-1} < t \leqslant 2^{-n}\)，\(\mathbf{f}(t)\) 是这样的序列：它的各项皆为零，除了指标为 n 的项等于 \(2^{n/p}\)（\(1 \leqslant p < +\infty\)）。试证 f 对 \(F^{\prime}\) 上的强拓扑是可测的，且属于 \(\Lambda_{\mathrm{F}'}^p\)，但任何阶梯函数序列 \((\mathbf{f}_n)\) 都不趋于 \(\mathbf{f}\)（在空间 \(\Lambda_{\mathrm{F}'}^p\) 中）。
\end{enumerate}

¶ 17) 设 \(\mu\) 是 \(T = [0,1]\) 上的 Lebesgue 测度，\(F\) 是 Banach 空间 \(L^1(\mathbf{N})\)，\(F' = L^\infty(\mathbf{N})\) 是其对偶。对每个 \(t \in [0,1[\)，设 \(t = \sum_{n=0}^\infty \xi_n 2^{-n}\) 为 \(t\) 的二进展开；用 \(f(t)\) 表示序列 \((\xi_n) \in F'\)，并令 \(f(1) = 0\)。

\begin{enumerate}
\def\labelenumi{\alph{enumi})}
\item
  试证：对 \(1 \leqslant p \leqslant +\infty\)，函数 \(\mathbf{f}\) 属于空间 \(\Lambda_{\mathrm{F}'}^p\)（参见习题 16）。
\item
  试证 \(\mathbf{f}\) 对拓扑 \(\sigma(\mathrm{F}',\mathrm{F})\) 是可测的，且对该拓扑存在一个阶梯函数序列几乎处处收敛于 \(\mathbf{f}\)。
\item
  试证 \(\mathbf{f}\) 对 \(\mathbf{F}'\) 上的强拓扑不是可测的（注意 \(|\mathbf{f}(t) - \mathbf{f}(t')| = 1\)，只要 \(0 \leqslant t < t' < 1\)）。
\item
  试证：在空间 \(\Lambda_{\mathrm{F}'}^p\)（\(1 \leqslant p < +\infty\)）中，\(\mathbf{f}\) 不是任何阶梯函数序列的极限。（可归结为只考虑这样的阶梯函数：它们是区间的特征函数的线性组合（系数取自 \(\mathbf{F}'\)），其中区间的端点形如 \(k/2^n\)；若 \(\mathbf{g}\) 是这样一个阶梯函数，试证在习题 16 b) 的记号下有 \(\mathrm{M}_1'(\mathbf{f} - \mathbf{g}) \geqslant 1/4\)。）
\item
  试证不存在任何映射 \(\mathbf{h}\)——\(\mathrm{T}\) 到 \(\mathrm{F}'\) 中的、对强拓扑可测的映射——使 \(\mathbf{f} - \mathbf{h}\) 是标量可忽略的。（注意对这样的映射，必然有 \(\mathrm{N}_{\infty}(\mathbf{h}) \leqslant 1\)；然后由此得出与 \(d\) 的结果矛盾。）
\end{enumerate}

\begin{enumerate}
\def\labelenumi{\arabic{enumi})}
\setcounter{enumi}{17}
\tightlist
\item
  试证：当把 \(\mathbf{f}\) 为 \(\mu\)-可测（对 \(\mathbf{F}\) 的原拓扑）这一条件换成 \(\mathbf{f}\) 为 \(\mu\)-可测（对弱化拓扑 \(\sigma(\mathrm{F},\mathrm{F}')\)）这一条件时，命题 8 仍然成立。（利用第 IV 章 §4 习题 19 c）。）
\end{enumerate}

¶ 19) 设 f 是 T 到 F 中的标量本质 μ-可积映射。对每个 \(z' \in F'\)，用 \(u_{\mathbf{f}}(\mathbf{z}')\) 表示 \(L^{1}(\mu)\) 中由函数 \(\langle f, z' \rangle\) 所确定的类。若对每个函数 \(g \in L^{\infty}(\mu)\) 都有 \(\int g f d\mu \in F\)，则称 f 是标量良可积的。

\begin{enumerate}
\def\labelenumi{\alph{enumi})}
\item
  为使 \(\mathbf{f}\) 是标量良可积的，必要且充分的是 \(u_{\mathbf{f}}\) 对拓扑 \(\sigma(\mathrm{F}',\mathrm{F})\) 与 \(\sigma(\mathrm{L}^1,\mathrm{L}^\infty)\) 连续。当此条件成立时，在 \(u_{\mathbf{f}}\) 之下，\(\mathrm{F}'\) 的每个等度连续子集的像是 \(\mathrm{L}^1\) 中的相对弱紧子集。特别地：\(\alpha\) 对每个等度连续子集 \(\mathrm{H}'\)（\(\mathrm{F}'\) 的）以及每个 \(\varepsilon >0\)，存在紧子集 \(\mathrm{L}\)（在 \(\mathrm{T}\) 中），使 \(\int_{\mathrm{T}-\mathrm{L}}|\langle\mathbf{f},\mathbf{z}'\rangle|d\mu\leqslant\varepsilon\) 对一切 \(\mathbf{z}'\in\mathrm{H}'\) 成立；\(\beta\) 对每个等度连续子集 \(\mathrm{H}'\)（\(\mathrm{F}'\) 的）以及每个 \(\varepsilon>0\)，存在 \(\eta>0\)，使得对每个开集 \(\mathrm{U}\subset\mathrm{T}\)（测度 \(\mu(\mathrm{U})\leqslant\eta\)），都有 \(\int_{\mathrm{U}}|\langle\mathbf{f},\mathbf{z}'\rangle|d\mu\leqslant\varepsilon\) 对一切 \(\mathbf{z}'\in\mathrm{H}'\) 成立（第 V 章 §5，习题 15）。
\item
  假设条件 \(\alpha\) 与 \(\beta\)（\(a\) 中的）都成立，此外还存在子空间 \(\mathrm{H} \subset \mathrm{L}^{\infty}\)，它对拓扑 \(\sigma(\mathrm{L}^{\infty}, \mathrm{L}^{1})\) 稠密，使得对每个函数 \(g \in \mathcal{L}^{\infty}\)，只要其类属于 \(\mathrm{H}\)，就有 \(\int g\mathbf{f}d\mu \in \mathrm{F}\)。试证在这些条件下，若再假设 \(\mathrm{F}\) 是拟完备的，则 \(\mathbf{f}\) 是标量良可积的。（先化归到 \(\mathrm{F}\)（对其原拓扑）完备的情形，并注意 \(\dot{g} \mapsto \int g\mathbf{f}d\mu\) 是 \(\mathrm{L}^{\infty}\) 到 \(\mathrm{F}'*\) 中的连续线性映射，其中 \(\mathrm{L}^{\infty}\) 赋以拓扑 \(\tau(\mathrm{L}^{\infty}, \mathrm{L}^{1})\)，\(\mathrm{F}'*\) 赋以在 \(\mathrm{F}'\) 的等度连续子集上一致收敛的拓扑。）
\item
  设 \(\mathbf{f}\) 是标量良可积的。试证对每个 Cauchy 序列 \((\mathbf{z}_n')\)（对拓扑 \(\sigma(\mathrm{F}',\mathrm{F})\)），序列 \((u_{\mathbf{f}}(\mathbf{z}_n'))\) 在 \(\mathrm{L}^1\) 中强收敛（利用第 IV 章 §5 习题 25 c) 与第 V 章 §5 习题 16 b)）。由此推出：若 A 是 \(\mathrm{L}^{\infty}\) 的单位球在映射 \(\dot{g} \mapsto \int g\mathbf{f}d\mu\) 下的像，则每个 Cauchy 序列 \((\mathbf{z}_n')\)（对拓扑 \(\sigma(\mathrm{F}',\mathrm{F})\)）都是对在 A 上一致收敛的拓扑的 Cauchy 序列。
\end{enumerate}

\(d)\) 由 \(c)\) 推出：若 \(\mathbf{f}\) 是标量良可积的，则 \(u_{\mathbf{f}}\) 把 \(\mathrm{F}'\) 的每个对 \(\sigma (\mathrm{F}',\mathrm{F})\) 等度连续且可度量化的子集映为 \(L^{1}\) 中的相对强紧集。特别地，若 F 是 Hausdorff 局部凸空间 G 的对偶 \(G'\)，F 赋以一个与 F 和 G 之间的对偶相容的拓扑，且 \(G'\) 中存在可数的强稠密集，则 \(u_{f}\) 把 \(F' = G\) 的每个有界子集映为 \(L^{1}\) 中相对强紧的集（把 G 嵌入其双对偶）。

\begin{enumerate}
\def\labelenumi{\arabic{enumi})}
\setcounter{enumi}{19}
\tightlist
\item
  设 F 是 Hausdorff 局部凸空间，\(F'\) 是其对偶，\(S'\) 是由这样的子集 \(A' \subset F'\) 组成的集：\(A'\) 的每个点序列都有对 \(\sigma(\mathrm{F}',\mathrm{F})\) 收敛的子序列。试证：对 F 的有界子集 A，下列条件等价：\(\alpha\) A 对 \(S'\)-拓扑是预紧的；\(\beta\) 每个对 \(\sigma(\mathrm{F}',\mathrm{F})\) 收敛的序列都在 A 上一致收敛。（设 S 是由 F 与 A 的有限子集组成的集；注意 \(\beta\) 蕴含每个集 \(A' \in S'\) 对 S-拓扑是预紧的，这里要用 GT，II，§4 习题 6。然后应用 TVS，IV，§1 习题 4。）
\end{enumerate}

¶ 21) 设 f 是 T 到拟完备 Hausdorff 局部凸空间 F 中的标量良可积映射。

\begin{enumerate}
\def\labelenumi{\alph{enumi})}
\item
  试证在下列每种情形中，映射 \(\dot{g} \mapsto \int g\mathbf{f}  d\mu\) 把 \(\mathrm{L}^{\infty}\) 的单位球映为 \(\mathbf{F}\) 中对原拓扑的相对紧子集：\(1^{\circ}\) \(\mathbf{F}\) 中存在可数稠密集；\(2^{\circ}\) \(\mathbf{F}\) 是一个对偶 \(\mathbf{G}'\)——其中 \(\mathbf{G}\) 或者可度量化，或者是可度量化空间序列的严格直接极限——且 \(\mathbf{F}\) 赋以拓扑 \(\tau (\mathbf{G}',\mathbf{G})\)。（利用习题 19 c) 与习题 20；在情形 \(2^{\circ}\)，用 Šmulian 定理（TVS，IV，§5，习题 2 c））。）
\item
  试证 a) 的结论在下述情况下仍然成立：不对 F 作新的假设，而假设 f 是可测的。（把 F 嵌入一族 Banach 空间的乘积，化归到 F 是 Banach 空间的情形；然后利用习题 19 a)，化归到 T 是紧的情形；再应用 a) 的情形 \(1^{\circ}\)。
\end{enumerate}

\begin{enumerate}
\def\labelenumi{\arabic{enumi})}
\setcounter{enumi}{21}
\tightlist
\item
  在 Hausdorff 局部凸空间 \(\mathbf{F}\) 中，设 \((\mathbf{z}_{\alpha})_{\alpha \in \mathbf{A}}\) 是这样的族：映射 \(\alpha \mapsto \mathbf{z}_{\alpha}\) 当 \(\mathbf{A}\) 赋以由每点质量 \(+1\) 定义的测度时是标量良可积的（习题 19）。
\end{enumerate}

\begin{enumerate}
\def\labelenumi{\alph{enumi})}
\item
  试证族 \((\mathbf{z}_{\alpha})\) 对每个与 F 和 \(F'\) 之间的对偶相容的拓扑 T 都是可和的（利用习题 19 a)）。当 F 对 T 拟完备时，逆命题成立。
\item
  设 \(F = G'\)，\(F' = G\)，其中 G 是这样的 Hausdorff 局部凸空间：\(G'\) 中存在可数的强稠密集。试证 \((\mathbf{z}_{\alpha})\) 对 \(G'\) 上的强拓扑是可和的（该拓扑不一定与 \(G'\) 与 G 之间的对偶相容（参见习题 19d））。试证这一结果不能推广到 \(G = L^{1}(\mathbf{N})\)、\(G' = L^{\infty}(\mathbf{N})\) 的情形。
\end{enumerate}

\begin{enumerate}
\def\labelenumi{\arabic{enumi})}
\setcounter{enumi}{22}
\tightlist
\item
  设 \(\mu\) 是 \(T\) 上的正测度，由可数个紧集的并所承载。
\end{enumerate}

\begin{enumerate}
\def\labelenumi{\alph{enumi})}
\tightlist
\item
  设 H 是由 \(\mu\)-可测数值函数组成的集，使得对每个 \(t \in \mathbf{T}\)，\(\sup_{f \in \mathrm{H}} f(t) < +\infty\)。试证存在一个有限值的 \(\mu\)-可测函数 \(g\)，使得
\end{enumerate}

\(f(t) \leqslant g(t)\) 对每个函数 \(f \in \mathrm{H}\) 都几乎处处成立。（化归到 \(\mu\) 有界且 \(\mathbf{H}\) 中函数取值于 \([-1, +1]\) 的情形，这可借助 \(\overline{\mathbf{R}}\) 到该区间上的一个递增同胚做到。然后考虑 \(\widetilde{g}\)，即函数的类所成之集在 \(\mathrm{L}^1(\mu)\) 中的上确界，其中函数取自 \(\mathbf{H}\)（第 IV 章 §3，第 6 号，命题 14），并注意存在 \(\mathbf{H}\) 中的函数序列几乎处处收敛于 \(g\)。）

\begin{enumerate}
\def\labelenumi{\alph{enumi})}
\setcounter{enumi}{1}
\item
  设 \(\mathbf{f}\) 是一个标量 \(\mu\)-可测映射，定义在 \(\mathrm{T}\) 上、取值于 Hausdorff 局部凸空间 \(\mathrm{F}\)。对每个子集 \(\mathrm{B}'\)（\(\mathrm{F}'\) 的、对 \(\sigma(\mathrm{F}',\mathrm{F})\) 有界的），试证存在可测的有限值数值函数 \(g_{\mathrm{B}'}\)，使得对每个 \(\mathbf{z}' \in \mathrm{B}'\)，\(|\langle \mathbf{f}(t),\mathbf{z}' \rangle| \leqslant g_{\mathrm{B}'}(t)\) 几乎处处成立。
\item
  设 \(\mathbf{F}\) 是 Fréchet 空间，\(\mathbf{f}\) 是一个标量 \(\mu\)-可测映射，定义在 \(\mathbf{T}\) 上、取值于 \(\mathbf{F}\)。试证：对每个紧子集 \(\mathbf{K}\)（\(\mathbf{T}\) 的）以及每个 \(\varepsilon > 0\)，存在紧集 \(\mathrm{K}' \subset \mathrm{K}\)，使 \(\mu(\mathrm{K} - \mathrm{K}') \leqslant \varepsilon\) 且 \(\mathbf{f}\) 在 \(\mathrm{K}'\) 上的限制有界（利用 \(b\)））。
\end{enumerate}

\begin{enumerate}
\def\labelenumi{\arabic{enumi})}
\setcounter{enumi}{23}
\item
  设 F 是完备的 Hausdorff 局部凸空间，包含可数稠密集。设 f 是 T 到 F 中的映射，标量 \(\mu\)-可测且标量本性有界；试证：对每个函数 \(g \in L^{1}\)，gf 是标量可积的且 \(\int g f d\mu \in F\)。（利用 TVS，III，§3，第 6 号定理 2 的推论 1，注意 \(F'\) 的等度连续子集对 \(\sigma(F', F)\) 可度量化，并应用 Lebesgue 定理。）
\item
  设 F 是可分 Fréchet 空间，其中每个对 \(\sigma(F, F')\) 的 Cauchy 序列都对这一拓扑收敛（例如空间 \(L^1(T', \nu)\)，其中 \(T'\) 具有可数基（第 V 章 §5，习题 16 b）））。试证 T 到 F 中的每个标量本质可积映射 f 都是标量良可积的。（先证明对每个具有紧支集的函数 \(g \in \mathcal{L}^\infty(\mu)\)，有 \(\int g\mathbf{f} d\mu \in F\)，这要注意到 f 是 \(\mu\)-可测的（第 5 号，命题 12），并应用第 2 号的命题 8 与假设。然后借助假设由此推出：\(\int_A g\mathbf{f} d\mu \in F\) 对每个函数 \(g \in \mathcal{L}^\infty\) 以及每个为可数个紧集之并的集 A 都成立。最后，用习题 19 的记号，证明在 \(u_{\mathbf{f}}\) 之下，\(F'\) 的每个等度连续子集的像都在 \(L^1\) 中相对弱紧，这里要用 Eberlein 定理（TVS，IV，§5，第 3 号，定理 1）以及 \(F'\) 的等度连续子集对 \(\sigma(F', F)\) 可度量化这一事实。）
\item
  设 \((f_n)\) 是 \(\mu\)-可积函数序列，定义在 \(\mathbf{T}\) 上，使得对几乎所有 \(t \in \mathbf{T}\)，\(\sup_{n} |f_n(t)| < +\infty\) 且 \(\sup_{n} \int |f_n(t)| d\mu(t) = +\infty\)。试证存在标量序列 \((c_n)\)，使 \(\sum_{n} |c_n| < +\infty\) 且 \(\sum_{n} c_n f_n\) 不是 \(\mu\)-可积的（把 Gelfand-Dunford 定理应用于空间 \(\mathbf{L}^1(\mathbf{N}) = \mathbf{F}\)。）
\item
  设 E 是可分实 Banach 空间，\(G = E'\) 是其对偶，\(G'\) 是 E 的子空间，在 E 中稠密、异于 E 且是桶式的（TVS，III，§4，习题 4）。试证：若 G 赋以拓扑 \(\sigma(G, G')\)，其对偶 \(G' = \mathcal{L}(G; \mathbf{R})\) 赋以紧收敛拓扑，则 \(G'\) 不是拟完备的，且存在序列 \((x_n)\) 在 \(G'\) 中趋于 0，以及一个数序列 \((\lambda_n)\)——这些数 \textgreater0 且和有限——使得通项为 \(\lambda_n x_n\) 的级数在 \(G'\) 中不收敛。（先注意 G 的每个对 \(\sigma(G, G')\) 有界的子集对 G 上的范数有界。取一点 \(a \in E\) 不属于 \(G'\)，再取点序列 \((a_n)\)（\(G'\) 中的点），它在 E 中收敛于 a 且满足 \(\|a_{n+1} - a_n\| \leqslant 1/2^n\)。）
\end{enumerate}

\setcounter{exercisegroup}{2}
\edef\CurrentExerciseGroup{\arabic{exercisegroup}}
